% =====================================================================
%  Arabic-language abstract (الملخّص).
%  Typeset with native Arabic (polyglossia + Amiri), set up in preamble.tex.
%  REQUIRES a Unicode engine (XeLaTeX/LuaLaTeX). Under pdfLaTeX it degrades
%  to a one-line notice so the document still compiles.
%  Latin technical tokens are wrapped in \textenglish{...} so they render
%  left-to-right in Latin Modern; numerals=maghrib keeps digits Western.
% =====================================================================

\ifPDFTeX
\chapter*{Abstract (Arabic)}
\addcontentsline{toc}{chapter}{Abstract (Arabic)}
\noindent\textit{The Arabic-language abstract is typeset with native Arabic and
requires compilation with XeLaTeX or LuaLaTeX (\texttt{latexmk main.tex}).}
\else
\begin{otherlanguage}{arabic}
\chapter*{المستخلص}
\addcontentsline{toc}{chapter}{\textenglish{Abstract (Arabic)}}

توثّق هذه الرسالة تصميم نظام تايم لنز (\textenglish{\TimeLens{}}) وتنفيذه وتقييمه، وهو دليلٌ
إرشاديٌّ للأجهزة المحمولة مدعومٌ بالذكاء الاصطناعي والواقع المعزَّز، ومخصَّصٌ
للمتحف المصري الكبير (\textenglish{GEM}). فبمجرّد توجيه الهاتف نحو إحدى القطع
المعروضة، يرى الزائرُ القطعةَ الأثرية وقد جرى التعرّف عليها في الزمن الحقيقي،
ويمكنه طرح أسئلة متابعة تُجاب بالعربية أو الإنجليزية. ويتصدّى هذا العمل لثلاث
مشكلات ترتبط تحديدًا بالتشغيل داخل قاعات العرض: التشابه البصري الدقيق بين
\textenglish{51} قطعة أثرية مُفهرَسة (كثيرٌ منها تماثيل رمسيسية شبه متطابقة)، والفجوة بين بيانات
التدريب المنسَّقة وظروف التصوير بكاميرا الهاتف المحمول، وخطر أن يذكر الدليلُ
الذكيُّ «حقائق» تاريخية لا يستطيع إسنادها.

تُقدّم الرسالة إسهامَين بحثيَّين هندسيَّين ودراسةً سوقية. الإسهام الأول كاشفٌ
للقطع الأثرية يعمل على الجهاز نفسه، طُوِّر عبر دراسة تَكرارية موجَّهة بجودة
البيانات: بدءًا من التوسيم الآلي بنموذج أساس (\textenglish{YOLO-World})، مرورًا
بقواعد مكانية لتنقية التوسيمات، وصولًا إلى مجموعة بيانات موسومة يدويًا بالكامل.
وتَعزِل الدراسةُ جودةَ التوسيم — لا معماريةَ النموذج ولا دقةَ الدخل — بوصفها
العاملَ الحاسم: إذ يَحُلّ نموذج \textenglish{YOLOv8n} النهائي كلَّ فئة كانت تفشل
سابقًا، مع بقائه أصلًا بصيغة \textenglish{TensorFlow Lite} لا يتجاوز حجمه
\textenglish{5.97} ميغابايت ويعمل في الزمن الحقيقي على هاتف متوسط الإمكانات. أمّا الإسهام
الثاني فهو دليلٌ ثنائي اللغة قائمٌ على التوليد المعزَّز بالاسترجاع
(\textenglish{RAG}) ومُسنَدٌ إلى قاعدة معرفة من 108 سجلات في \textenglish{ChromaDB}؛
وقد اختار تقييمٌ منهجيٌّ لسبعة نماذج لغوية مرشَّحة النموذجَ
\textenglish{Gemma 4 E2B (Q4\_K\_M)}، كما خفّض تحسينٌ من عشر خطوات زمنَ الاستجابة
الكلي من أكثر من \textenglish{30} ثانية إلى نحو \textenglish{10} ثوانٍ. ويُدمَج الإسهامان معًا في تطبيق
\textenglish{Flutter} جاهزٍ للإنتاج، مع واجهة ثنائية اللغة، وتقييدٍ حسب موقع
المتحف، وخاصيةِ تحويل النص إلى كلام.

% وإلى جانب النتائج الرئيسية، كُتبت الرسالة لتكون قابلةً لإعادة الإنتاج وصادقةً
% بشأن حدودها: إذ تُفصِح دراسة الكشف صراحةً عن تسرّب الإطارات في بيانات التحقّق
% وعن انزياح المجال بين بيئة التدريب وبيئة التشغيل، وتُورِد نتيجة متحفِّظة على
% بيانات محجوزة (\textenglish{mAP@0.5 = 0.751}) إلى جانب نتائج البيانات داخل
% التوزيع. ويتضمّن المستند ملاحظاتٍ لإعادة الإنتاج، وبيانًا للأخلاقيات والتحيُّز،
% وإفصاحًا رسميًا عن استخدام الذكاء الاصطناعي التوليدي.

وإلى جانب المقاييس الرئيسية، تُورِد الرسالة نتيجةً متحفِّظة على بيانات
محجوزة (\textenglish{mAP@0.5 = 0.751}) إلى جانب نتائج البيانات داخل
التوزيع. ويتضمّن المستند ملاحظاتٍ لإعادة الإنتاج، وبيانًا للأخلاقيات
والتحيُّز، وإفصاحًا رسميًا عن استخدام الذكاء الاصطناعي التوليدي.

\vspace{1.2em}
\noindent\textbf{الكلمات المفتاحية:} كشف الأجسام على الجهاز؛ \textenglish{YOLOv8n}؛
\textenglish{TensorFlow Lite}؛ التوليد المعزَّز بالاسترجاع (\textenglish{RAG})؛
\textenglish{Gemma}؛ \textenglish{ChromaDB}؛ الإجابة عن الأسئلة ثنائية اللغة؛
الواقع المعزَّز؛ \textenglish{Flutter}؛ حوسبة التراث الثقافي؛ المتحف المصري الكبير.
\end{otherlanguage}
\fi
